% Options for packages loaded elsewhere
\PassOptionsToPackage{unicode}{hyperref}
\PassOptionsToPackage{hyphens}{url}
\documentclass[
]{article}
\usepackage{xcolor}
\usepackage[letterpaper,top=2cm,bottom=2cm,left=1.5cm,right=1.5cm]{geometry}
\usepackage{amsmath,amssymb}
\setcounter{secnumdepth}{-\maxdimen} % remove section numbering
\usepackage{iftex}
\ifPDFTeX
  \usepackage[T1]{fontenc}
  \usepackage[utf8]{inputenc}
  \usepackage{textcomp} % provide euro and other symbols
\else % if luatex or xetex
  \usepackage{unicode-math} % this also loads fontspec
  \defaultfontfeatures{Scale=MatchLowercase}
  \defaultfontfeatures[\rmfamily]{Ligatures=TeX,Scale=1}
\fi
\usepackage{lmodern}
\ifPDFTeX\else
  % xetex/luatex font selection
\fi
% Use upquote if available, for straight quotes in verbatim environments
\IfFileExists{upquote.sty}{\usepackage{upquote}}{}
\IfFileExists{microtype.sty}{% use microtype if available
  \usepackage[]{microtype}
  \UseMicrotypeSet[protrusion]{basicmath} % disable protrusion for tt fonts
}{}
\makeatletter
\@ifundefined{KOMAClassName}{% if non-KOMA class
  \IfFileExists{parskip.sty}{%
    \usepackage{parskip}
  }{% else
    \setlength{\parindent}{0pt}
    \setlength{\parskip}{6pt plus 2pt minus 1pt}}
}{% if KOMA class
  \KOMAoptions{parskip=half}}
\makeatother
\usepackage{longtable,booktabs,array}
\usepackage{caption}
\captionsetup[table]{skip=6pt}
\newcounter{none} % for unnumbered tables
\usepackage{calc} % for calculating minipage widths
% Correct order of tables after \paragraph or \subparagraph
\usepackage{etoolbox}
\makeatletter
\patchcmd\longtable{\par}{\if@noskipsec\mbox{}\fi\par}{}{}
\makeatother
% Allow footnotes in longtable head/foot
\IfFileExists{footnotehyper.sty}{\usepackage{footnotehyper}}{\usepackage{footnote}}
\makesavenoteenv{longtable}
\setlength{\emergencystretch}{3em} % prevent overfull lines
\providecommand{\tightlist}{%
  \setlength{\itemsep}{0pt}\setlength{\parskip}{0pt}}
\usepackage{bookmark}
\IfFileExists{xurl.sty}{\usepackage{xurl}}{} % add URL line breaks if available
\urlstyle{same}
% fallback for those not using the hyperref driver hyperxmp:
\makeatletter
\@ifundefined{xmpquote}{\newcommand{\xmpquote}[1]{#1}}{}
\makeatother
\hypersetup{
  hidelinks,
  pdfcreator={LaTeX via pandoc}}

\author{}
\date{}

\begin{document}

\textbf{ Executive Report }

Factor Analysis vs Principal Component Analysis

Recovering a Simulated Construct Structure: A Worked Example

\begin{center}\rule{0.5\linewidth}{0.5pt}\end{center}

Analysis of 200 Students Across Nine Assessment Variables

\section{Executive Summary}

\textbf{Synthetic data notice:} This report analyzes a single
\textbf{simulated} sample (\texttt{fetch\_educational.py}) generated
from a known three-factor model with factor correlations of 0.20-0.30.
The ``students'' and their scores are not real. The purpose of this
report is to demonstrate whether Factor Analysis and PCA can
\textbf{recover a structure we already know}, not to validate a real
assessment instrument. Genuine construct validation requires independent
real-world data, replication, and typically a confirmatory factor
analysis (CFA).

This report presents a worked example examining whether nine simulated
educational assessment variables recover three distinct underlying
constructs: Quantitative Reasoning, Verbal Ability, and Interpersonal
Skills. The analysis employs Factor Analysis (FA) with \textbf{Promax}
(oblique) rotation as the primary method -- because the data were
simulated with correlated factors, an oblique rotation is the
appropriate default -- with Varimax (orthogonal) rotation fit alongside
it as a sensitivity comparison, plus Principal Component Analysis (PCA)
on data from 200 simulated students.

\textbf{Key Finding:} Factor Analysis with Promax rotation recovers the
three simulated constructs with high clarity: each assessment loads
strongly (\textgreater0.87) on exactly one factor, and the recovered
factor correlations (Φ ≈ 0.19-0.36) are consistent with the 0.20-0.30
correlations used to generate the data.

\textbf{Recommendation:} This worked example illustrates the mechanics
of FA/PCA and rotation choice. It does not, by itself, validate any real
assessment battery; that would require independent data and further
evidence (e.g. CFA, reliability, criterion validity).

\section{Research Question}

\textbf{Can Factor Analysis and PCA recover the three simulated
underlying constructs (Quantitative, Verbal, Interpersonal) in this
synthetic dataset, and how does the choice between oblique and
orthogonal rotation affect the answer?}

This worked example illustrates:

\begin{itemize}
\item
  How rotation choice (oblique vs. orthogonal) interacts with a
  data-generating process that has correlated factors
\item
  How to read loadings, communalities, and a factor-correlation matrix
\item
  The difference between a method \textbf{recovering a known structure}
  and \textbf{validating a real construct}
\end{itemize}

\section{Dataset Description}

\textbf{Sample:} 200 simulated students (\texttt{fetch\_educational.py},
synthetic data)

\textbf{Variables:} Nine assessment scores across three theoretical
domains:

\begin{itemize}
\item
  \textbf{Quantitative Domain:} MathScore, AlgebraScore, GeometryScore
\item
  \textbf{Verbal Domain:} ReadingComp, Vocabulary, Writing
\item
  \textbf{Interpersonal Domain:} Collaboration, Leadership,
  Communication
\end{itemize}

\textbf{Data-generating process:} The three latent factors were
simulated as \textbf{correlated} (0.20-0.30), not independent -- this is
the key fact that motivates using an oblique rotation as the primary
analysis below.

All variables were standardized (mean=0, std=1) prior to analysis to
ensure equal contribution regardless of original measurement scales.

\begin{center}\rule{0.5\linewidth}{0.5pt}\end{center}

\section{Methodological Approach}

\subsection{Factor Analysis}

Factor Analysis (FA) is a statistical technique that identifies latent
constructs underlying observed variables. Unlike PCA, FA distinguishes
between:

\begin{itemize}
\item
  \textbf{Common variance:} Shared among variables, explained by factors
\item
  \textbf{Unique variance:} Specific to each variable, combining
  variable-specific true variance and measurement error (these two
  cannot be separated from uniqueness alone)
\end{itemize}

\textbf{Key Parameters:}

\begin{itemize}
\item
  Extraction method: Principal Axis Factoring (PAF)
\item
  Number of factors: 3 (theoretical expectation)
\item
  Rotation: \textbf{Promax} (oblique, primary analysis -- allows
  correlated factors, matching the data-generating process), with
  \textbf{Varimax} (orthogonal) fit as a sensitivity comparison
\end{itemize}

\subsection{Principal Component Analysis}

PCA identifies linear combinations of variables that capture maximum
variance. It does not distinguish common from unique variance, using all
available variance.

\textbf{Key Parameters:}

\begin{itemize}
\item
  All components extracted (9 total)
\item
  No rotation applied
\item
  Focus on first 3 components for comparison
\end{itemize}

\subsection{Statistical Assumptions}

Two critical tests verified data suitability for Factor Analysis:

\textbf{Bartlett's Test of Sphericity}

\begin{itemize}
\item
  Chi-square = 1393.938, p \textless{} 0.001
\item
  Confirms variables are sufficiently correlated for FA
\item
  Rejects hypothesis that correlation matrix is an identity matrix
\end{itemize}

\textbf{Kaiser-Meyer-Olkin (KMO) Measure}

\begin{itemize}
\item
  Overall KMO = 0.799 (Acceptable)
\item
  Individual MSA values: 0.728 to 0.850 (all adequate)
\item
  Confirms sampling adequacy for factor extraction
\end{itemize}

\textbf{Conclusion:} Both tests passed. Data meets requirements for
meaningful Factor Analysis.

\begin{center}\rule{0.5\linewidth}{0.5pt}\end{center}

\section{Factor Analysis Results}

\subsection{Eigenvalues and Variance Explained}

Three factors extracted with eigenvalues exceeding the Kaiser criterion
(\textgreater1.0):

{\def\LTcaptype{none} % do not increment counter
\begin{longtable}[]{@{}ccc@{}}
\toprule\noalign{}
\endhead
\bottomrule\noalign{}
\endlastfoot
\textbf{Factor} & \textbf{Eigenvalue} & \textbf{Cumulative Variance} \\
1 & 4.041 & 44.9\% \\
2 & 2.124 & 68.5\% \\
3 & 1.635 & 86.7\% \\
\end{longtable}
}

The three retained factors' communalities sum to 86.7\% of the
\textbf{total standardized variance} (sum of communalities / 9
variables) -- not ``86.7\% of common variance,'' which would be a
different (and here, trivially 100\%) quantity. Remaining factors have
eigenvalues \textless1.0, indicating they capture primarily noise.

\subsection{Communalities}

All variables showed high communalities (h²), indicating they are
well-explained by the three-factor model:

{\def\LTcaptype{none} % do not increment counter
\begin{longtable}[]{@{}lll@{}}
\toprule\noalign{}
\endhead
\bottomrule\noalign{}
\endlastfoot
\textbf{Variable} & \textbf{Communality (h²)} & \textbf{Uniqueness
(u²)} \\
MathScore & 0.879 & 0.121 \\
AlgebraScore & 0.905 & 0.095 \\
GeometryScore & 0.848 & 0.152 \\
ReadingComp & 0.852 & 0.148 \\
Vocabulary & 0.838 & 0.162 \\
Writing & 0.852 & 0.148 \\
Collaboration & 0.874 & 0.126 \\
Leadership & 0.873 & 0.127 \\
Communication & 0.880 & 0.120 \\
\end{longtable}
}

\textbf{Average communality = 0.867}: on average, 86.7\% of each
variable's variance is shared with the other variables through the three
common factors, and 13.3\% is unique to that variable (specific variance
plus measurement error, not separable from communalities alone).

\begin{center}\rule{0.5\linewidth}{0.5pt}\end{center}

\subsection{Unrotated Factor Interpretation}

The unrotated solution revealed three problematic patterns:

\textbf{Factor 1: General Educational Ability}

\begin{itemize}
\item
  All nine variables load positively (0.577 to 0.737)
\item
  Represents overall academic competence, not specific abilities
\item
  Problem: Cannot distinguish which skills a student possesses
\end{itemize}

\textbf{Factor 2: Interpersonal vs. Cognitive Contrast}

\begin{itemize}
\item
  Positive loadings: Collaboration (0.706), Leadership (0.710),
  Communication (0.650)
\item
  Negative loadings: Math (-0.478), Algebra (-0.438), Geometry (-0.464)
\item
  Problem: Bipolar structure creates ambiguous interpretation
\end{itemize}

\textbf{Factor 3: Quantitative vs. Verbal Contrast}

\begin{itemize}
\item
  Positive loadings: Math (0.422), Algebra (0.555), Geometry (0.338)
\item
  Negative loadings: Reading (-0.527), Vocabulary (-0.588), Writing
  (-0.558)
\item
  Problem: Suggests students cannot excel at both domains
\end{itemize}

\textbf{Critical Issues:}

\begin{enumerate}
\item
  Lack of simple structure (variables load on multiple factors)
\item
  Bipolar factors with confusing negative loadings
\item
  Not theory-aligned (expected three positive constructs, got contrasts
  + general factor)
\item
  Low practical utility for subscale creation
\end{enumerate}

\textbf{Explanation:} Unrotated factors maximize variance extraction,
not interpretability. This is why rotation is essential.

\begin{center}\rule{0.5\linewidth}{0.5pt}\end{center}

\subsection{Rotated Factor Solution (Promax, Primary)}

Promax rotation transformed factors to achieve simple structure while
allowing them to correlate:

{\def\LTcaptype{none} % do not increment counter
\begin{longtable}[]{@{}cccc@{}}
\toprule\noalign{}
\endhead
\bottomrule\noalign{}
\endlastfoot
\textbf{Variable} & \textbf{Factor 1} & \textbf{Factor 2} &
\textbf{Factor 3} \\
MathScore & \textbf{0.931} & -0.018 & 0.027 \\
AlgebraScore & \textbf{0.984} & 0.036 & -0.134 \\
GeometryScore & \textbf{0.874} & -0.018 & 0.121 \\
ReadingComp & 0.051 & -0.011 & \textbf{0.906} \\
Vocabulary & -0.032 & -0.032 & \textbf{0.934} \\
Writing & -0.023 & 0.047 & \textbf{0.917} \\
Collaboration & 0.001 & \textbf{0.945} & -0.041 \\
Leadership & -0.001 & \textbf{0.947} & -0.052 \\
Communication & 0.002 & \textbf{0.902} & 0.110 \\
\end{longtable}
}

\textbf{Bold values} indicate salient loadings (\textgreater0.4
threshold).

\textbf{Factor Correlation Matrix (Φ):} Unlike Varimax, Promax does not
force factors to be uncorrelated:

{\def\LTcaptype{none} % do not increment counter
\begin{longtable}[]{@{}cccc@{}}
\toprule\noalign{}
\endhead
\bottomrule\noalign{}
\endlastfoot
& \textbf{Factor 1} & \textbf{Factor 2} & \textbf{Factor 3} \\
\textbf{Factor 1} & 1.000 & 0.194 & 0.357 \\
\textbf{Factor 2} & 0.194 & 1.000 & 0.271 \\
\textbf{Factor 3} & 0.357 & 0.271 & 1.000 \\
\end{longtable}
}

These recovered correlations (0.19-0.36) are close to the 0.20-0.30
correlations actually used to simulate the data -- evidence that Promax
is recovering the true generating structure, whereas Varimax's
assumption of zero factor correlation would misrepresent it.

\textbf{Interpretation:}

\begin{itemize}
\item
  \textbf{Factor 1 = Quantitative Reasoning:} Math, Algebra, Geometry
  (loadings \textgreater0.87)
\item
  \textbf{Factor 2 = Interpersonal Skills:} Collaboration, Leadership,
  Communication (loadings \textgreater0.90)
\item
  \textbf{Factor 3 = Verbal Ability:} Reading, Vocabulary, Writing
  (loadings \textgreater0.89)
\end{itemize}

\textbf{Result:} Each variable loads strongly on exactly one factor,
with small cross-loadings. In this synthetic sample, this recovers the
theoretical three-construct model used to generate the data -- it
demonstrates the method, and is not itself evidence that a \textbf{real}
assessment battery would show the same pattern.

\subsection{Sensitivity Comparison: Varimax (Orthogonal)}

{\def\LTcaptype{none} % do not increment counter
\begin{longtable}[]{@{}cccc@{}}
\toprule\noalign{}
\endhead
\bottomrule\noalign{}
\endlastfoot
\textbf{Variable} & \textbf{Factor 1} & \textbf{Factor 2} &
\textbf{Factor 3} \\
MathScore & \textbf{0.916} & 0.067 & 0.188 \\
AlgebraScore & \textbf{0.944} & 0.105 & 0.048 \\
GeometryScore & \textbf{0.877} & 0.075 & 0.270 \\
ReadingComp & 0.210 & 0.113 & \textbf{0.892} \\
Vocabulary & 0.131 & 0.089 & \textbf{0.901} \\
Writing & 0.144 & 0.166 & \textbf{0.897} \\
Collaboration & 0.075 & \textbf{0.928} & 0.081 \\
Leadership & 0.072 & \textbf{0.929} & 0.071 \\
Communication & 0.099 & \textbf{0.906} & 0.223 \\
\end{longtable}
}

Varimax gives similar loadings and the same variable groupings here, but
by construction reports zero correlation between factors (Φ = I), which
does not match how this dataset was actually generated. We report it
only as a sensitivity check, not as the primary result.

\begin{center}\rule{0.5\linewidth}{0.5pt}\end{center}

\section{Principal Component Analysis Results}

\subsection{Eigenvalues and Variance}

PCA extracted nine components with the following eigenvalue structure:

{\def\LTcaptype{none} % do not increment counter
\begin{longtable}[]{@{}ccc@{}}
\toprule\noalign{}
\endhead
\bottomrule\noalign{}
\endlastfoot
\textbf{Component} & \textbf{Eigenvalue} & \textbf{Cumulative
Variance} \\
1 & 4.062 & 44.9\% \\
2 & 2.135 & 68.5\% \\
3 & 1.643 & 86.7\% \\
4-9 & \textless1.0 & 13.3\% \\
\end{longtable}
}

PCA eigenvalues closely match FA eigenvalues for the first three
components, confirming three-dimensional structure. Components 4-9 have
eigenvalues \textless1.0 (Kaiser criterion), indicating they capture
mostly non-shared (unique + residual) variance -- not necessarily
``measurement noise'' specifically.

\subsection{Component Loadings}

Unlike FA's rotated solution, PCA loadings (unrotated) show more
distributed patterns:

\textbf{PC1 (44.9\% variance):} All variables load moderately
(0.28-0.37), suggesting a general factor

\textbf{PC2 (23.6\% variance):} Contrasts interpersonal skills
(positive) vs. cognitive skills (negative)

\textbf{PC3 (18.2\% variance):} Contrasts verbal (positive) vs.
quantitative (negative)

\textbf{Observation:} Without rotation, PCA components are harder to
interpret as distinct constructs. This demonstrates why rotation is
useful for interpretability -- it does not, by itself, demonstrate
construct validity.

\begin{center}\rule{0.5\linewidth}{0.5pt}\end{center}

\section{Comparative Analysis: FA vs PCA}

{\def\LTcaptype{none} % do not increment counter
\begin{longtable}[]{@{}lll@{}}
\toprule\noalign{}
\endhead
\bottomrule\noalign{}
\endlastfoot
\textbf{Criterion} & \textbf{Factor Analysis} & \textbf{Principal
Component Analysis} \\
Variance explained & 86.7\% of total standardized variance, via common
factors & 86.7\% of total variance, via first 3 components \\
Rotation applied & Yes (Promax primary, Varimax comparison) & No \\
Factor/component correlation & Nonzero (Φ ≈ 0.19-0.36, Promax) & Zero by
construction \\
Simple structure & Clear (each variable loads mainly on one factor) &
Absent (distributed loadings) \\
Interpretability & High (clear factor labels in this simulation) &
Moderate (requires interpretation) \\
Recovers known simulated structure & Yes, closely & Dimensionality yes;
construct labels less clear \\
Practical utility & High (enables subscale creation), pending real-data
validation & Moderate (less clear assignments) \\
\end{longtable}
}

\textbf{Key Insight:} Both methods identify three-dimensional structure.
FA with rotation provides better interpretability here, and Promax
additionally recovers the known factor correlations that Varimax cannot
represent. This shows FA is well-suited to this \textbf{kind} of problem
(correlated latent constructs) -- it is not, on its own, ``construct
validation.''

\section{Visual Results}

Three visualizations were generated to support findings:

\textbf{1. Scree Plots (fa\_scree.png)}

\begin{itemize}
\item
  Clear ``elbow'' after third component/factor in both FA and PCA
\item
  Eigenvalues 4-9 below Kaiser criterion (1.0)
\item
  Visual confirmation of three-factor retention
\end{itemize}

\textbf{2. Loading Heatmaps (fa\_loadings.png)}

\begin{itemize}
\item
  Unrotated loadings show mixed patterns with moderate values
\item
  Rotated loadings show distinct blocks with strong values
\item
  Dramatic visual demonstration of rotation's impact on interpretability
\end{itemize}

\textbf{3. PCA Biplot (pca\_biplot.png)}

\begin{itemize}
\item
  Variable arrows cluster into three groups (Quantitative, Verbal,
  Interpersonal)
\item
  Student scores distributed across PC1 and PC2
\item
  Visual confirmation of three-domain structure
\end{itemize}

\begin{center}\rule{0.5\linewidth}{0.5pt}\end{center}

\section{Conclusions}

\subsection{Answer to Research Question}

\textbf{YES, in this synthetic sample, Factor Analysis recovers the
three simulated constructs.}

Factor Analysis with Promax rotation (primary analysis):

\begin{itemize}
\item
  Each assessment loads strongly (\textgreater0.87) on exactly one
  factor
\item
  Clear simple structure achieved
\item
  Factors align precisely with the constructs used to generate the data:

  \begin{itemize}
  \item
    Quantitative Reasoning (Math, Algebra, Geometry)
  \item
    Interpersonal Skills (Collaboration, Leadership, Communication)
  \item
    Verbal Ability (Reading, Vocabulary, Writing)
  \end{itemize}
\item
  Recovered factor correlations (Φ ≈ 0.19-0.36) are consistent with the
  0.20-0.30 correlations used to generate the data
\item
  High communalities (average = 0.867) mean most of each variable's
  variance is shared through the common factors -- this is not the same
  as ``minimal measurement error,'' since uniqueness bundles specific
  variance and error together
\item
  86.7\% of total standardized variance is attributed to the three
  common factors
\end{itemize}

This is a demonstration that the method works as expected on data with a
known structure. It is \textbf{not}, by itself, evidence about any real
assessment battery.

\subsection{What This Does and Does Not Establish}

This worked example shows that FA with an appropriate (oblique) rotation
can recover a known simulated structure, including its factor
correlations. It does \textbf{not} establish:

\begin{itemize}
\item
  That any real assessment battery has this structure
\item
  That the specific assessment items are valid measures of these
  constructs
\item
  That the constructs are ``separable and measurable'' in a real
  population
\end{itemize}

Establishing those claims for a real instrument requires independent
data, replication across samples, and typically a confirmatory factor
analysis (CFA) that tests this specific three-factor model against the
data rather than exploring for structure.

\subsection{Methodological Lessons}

\textbf{Rotation choice is a modeling decision, not a default}

\begin{itemize}
\item
  Unrotated factors maximized variance but lacked interpretability
\item
  Both Promax and Varimax achieve interpretable simple structure here,
  but only Promax's nonzero Φ matches how the data was actually
  generated
\item
  When factors could plausibly correlate (the usual case for
  psychological/educational constructs), an oblique rotation should be
  the default choice, with an orthogonal rotation reported only as a
  comparison
\end{itemize}

\textbf{FA vs PCA Selection}

\begin{itemize}
\item
  Both methods identified three dimensions
\item
  FA with rotation provided clearer factor interpretation and, via
  Promax, correctly represented factor correlation
\item
  FA's separation of common and unique variance is useful when a
  latent-variable model is the actual object of interest
\end{itemize}

\textbf{Assumption Testing}

\begin{itemize}
\item
  KMO and Bartlett's tests confirmed data suitability
\item
  All variables showed adequate sampling adequacy
\item
  Proper assumption testing prevents meaningless results, but passing
  them does not by itself validate substantive conclusions
\end{itemize}

\begin{center}\rule{0.5\linewidth}{0.5pt}\end{center}

\section{Recommendations}

\subsection{If This Structure Were Confirmed on Real Data}

These applications would follow \textbf{only after} independent,
real-data confirmation (e.g. CFA, replication) of a three-factor
structure -- they are not licensed by this synthetic-data demonstration
alone:

\textbf{1. Subscale Creation}

\begin{itemize}
\item
  Three subscale scores: Quantitative, Verbal, Interpersonal
\item
  Each subscale with three indicators with high loadings
\item
  Factor scores or simple sum scores for student reporting
\end{itemize}

\textbf{2. Diagnostic Assessment}

\begin{itemize}
\item
  Identify student strengths and weaknesses across three domains
\item
  Target interventions to specific ability areas rather than general
  ``academic support''
\item
  Track domain-specific progress over time
\end{itemize}

\textbf{3. Program Evaluation}

\begin{itemize}
\item
  Assess effectiveness of interventions targeting specific constructs
\item
  Determine whether programs improve targeted abilities without
  affecting others
\end{itemize}

\textbf{4. Research Applications}

\begin{itemize}
\item
  Motivate further construct-validity studies using real measures
\item
  Motivate hypotheses about domain-specific effects, to be tested on
  independent data
\end{itemize}

\subsection{Cautions and Limitations}

\textbf{This Is Synthetic Data}

\begin{itemize}
\item
  All 200 ``students'' and their scores were simulated from a known
  three-factor model
\item
  Recovering the known structure demonstrates the method; it is not
  itself evidence about any real population
\item
  None of the conclusions above should be applied to real assessment
  data without independent verification
\end{itemize}

\textbf{Sample Specificity (if applied to real data)}

\begin{itemize}
\item
  A single sample from one context is never sufficient for validation
\item
  Replication with different, independent samples is required
\item
  Factor structure may vary across populations
\end{itemize}

\textbf{Unique Variance}

\begin{itemize}
\item
  Despite high communalities (average = 0.867), 13.3\% of variance is
  unique to each variable
\item
  Unique variance combines measurement error and variable-specific true
  variance -- these cannot be separated from communalities/uniqueness
  alone
\item
  Reliability analysis (e.g. Cronbach's alpha) is needed to estimate
  measurement error specifically
\end{itemize}

\subsection{Future Directions}

\textbf{Confirmatory Factor Analysis}

\begin{itemize}
\item
  On real data, test this specific three-factor model using structural
  equation modeling
\item
  Evaluate model fit with chi-square, CFI, RMSEA indices
\item
  Compare alternative models (e.g., single-factor, hierarchical)
\end{itemize}

\textbf{Invariance Testing}

\begin{itemize}
\item
  Test whether factor structure holds across groups (gender, age,
  ethnicity)
\item
  Establish measurement equivalence before making group comparisons
\end{itemize}

\textbf{Rotation Sensitivity}

\begin{itemize}
\item
  This report already uses Promax (oblique) as primary, given the
  correlated factor structure
\item
  On real data, also compare other oblique rotations (e.g. Oblimin) to
  check robustness of the recovered Φ matrix
\end{itemize}

\textbf{Predictive Validity}

\begin{itemize}
\item
  On real data, examine whether subscales predict relevant outcomes
  (grades, career success)
\item
  Establish criterion-related validity
\item
  Test incremental validity of three separate scores vs. composite
\end{itemize}

\begin{center}\rule{0.5\linewidth}{0.5pt}\end{center}

Report prepared from statistical analysis conducted using Python
(factor\_analyzer, scikit-learn)

Dataset: Synthetic data, N=200 simulated students, p=9 assessment
variables

Methods: Factor Analysis (Principal Axis Factoring, Promax rotation
primary, Varimax comparison), Principal Component Analysis

\end{document}

